We now describe a polynomial-time separation algorithm for $\eqref{eq:RP-pos-cutting-planes}$ and then for $\eqref{eq:RP-neg}.$

Separation of \eqref{eq:RP-pos-cutting-planes} for fixed $i$ can be done in $\mathcal{O}(m  \log m)$ time using a greedy algorithm.
Without loss of generality, consider a fixed point $i$ and
\[
    2 \geq \gamma_1 \geq \gamma_2 \geq \cdots \geq \gamma_m \geq 0,
\]
each corresponding to a clause.
Assume the $z$-variables are ordered accordingly.
Let $(\bm{z}, \gamma)$ be a solution to the (partial) LP\@.
Note that $\bm{z} \geq 0$.
Observe that a violated constraint exists if and only if there exist
\[
    \tau_1 < \tau_2 < \cdots < \tau_t \in [m].
\] such that
\begin{equation}\label{eq:violated_pos_cons}
    \gamma > 2 - \sum_{\kappa = 1}^t \overline{\gamma}_{\tau_\kappa}z_{\tau_\kappa},
\end{equation}
where
\begin{align*}
    \overline\gamma_{\tau_1} &= 2 - \gamma_{\tau_1}, \\
    \overline\gamma_{\tau_j} &= 2 - \gamma_{\tau_j} - (2 - \gamma_{\tau_{j-1}}) = \gamma_{\tau_{j-1}} - \gamma_{\tau_j}, \quad j = 2, \dots, t.
\end{align*}
If the inequality holds for a minimizer of the right hand side, then we have found a new constraint.
If the inequality does not hold for this minimizer, we may conclude that no violated constraint exists.
Finding a minimizer of the right hand side is equivalent to finding a maximizer of
\begin{equation}\label{eq:sep-pos-sum}
    \sum_{\kappa = 1}^t \overline{\gamma}_{\tau_\kappa}z_{\tau_\kappa}.
\end{equation}

Denote $\gamma_0$ = 2 and $\tau_0 = 0$ for convenience.
We propose the following algorithm.

\begin{algorithm}[h!]
\caption{Algorithm to find a maximizer of \eqref{eq:sep-pos-sum}}
\label{alg:sep-pos}
\begin{algorithmic}[1]
\Require $\gamma, \bm{z}$
\State $i \gets 1$
\While{$\left\{\sigma \fw \gamma_\sigma < \gamma_{\tau_{i-1}}\right\} \neq \emptyset$}
    \State Choose a $\tau_i$ such that
    $z_{\tau_i} = \max\left\{z_\sigma : \gamma_{\sigma_i} < \gamma_{\tau_{i-1}}\right\}.$
    \State $i \gets i + 1.$
\EndWhile \\
\Return $(\tau_\kappa)_{\kappa=1,\dots,i-1}.$
\end{algorithmic}
\end{algorithm}

This algorithm can be implemented to run in $O(m\log m)$ time by first sorting (the indices of) $\bm{z}$
and iterating over the sorted list, keeping in mind that $\tau$ values must be increasing.

\begin{theorem}\label{thm:sep-pos}
    \autoref{alg:sep-pos} returns a subset \[\tau_1 < \tau_2 < \cdots < \tau_t \in [m]\] that maximizes
    \[
        \sum_{\kappa = 1}^t \overline{\gamma}_{\tau_\kappa}z_{\tau_\kappa}.
    \]
\end{theorem}

\begin{proof}
    Let
    \[
        \tau_1 < \tau_2 < \cdots < \tau_t \in [m]
    \]
    be a sequence.
    Define
    \[
        u_{\sigma} = \umax{\tau \geq \sigma} \; z_{\tau}, \tquad \sigma \in [m].
    \]
    Now observe that
    \begin{equation}\label{eq:sep-pos-proof-1}
        \begin{split}
            S(\tau) = \sum_{\kappa = 1}^t \overline{\gamma}_{\tau_\kappa}z_{\tau_\kappa}
            & = \sum_{\kappa = 1}^t z_{\tau_\kappa} \sum_{\sigma = \tau_{\kappa-1}+1}^{\tau_{\kappa}} (\gamma_{\sigma-1} - \gamma_{\sigma}) \\
            & = \sum_{\kappa = 1}^t \sum_{\sigma = \tau_{\kappa-1}+1}^{\tau_{\kappa}} z_{\tau_\kappa}(\gamma_{\sigma-1} - \gamma_{\sigma}) \\
            & \leq \sum_{\kappa = 1}^t \sum_{\sigma = \tau_{\kappa-1}+1}^{\tau_{\kappa}} u_{\sigma}(\gamma_{\sigma-1} - \gamma_{\sigma}) \\
            & \leq \sum_{\sigma = 1}^m u_{\sigma}(\gamma_{\sigma-1} - \gamma_\sigma).
        \end{split}
    \end{equation}
    The first inequality holds because
    \[
         \tau_\kappa \geq \sigma \quad \implies z_{\tau_\kappa} \leq u_\sigma,
    \]
    and $\gamma_{\sigma-1} - \gamma_\sigma \geq 0$ for all $\sigma \in [m]$.
    The second holds because the ranges $\{\tau_{\kappa-1}+1, \dots, \tau_\kappa\}, \; \kappa=1,\dots,t$
    partition $\{1, \dots, \tau_t\}$ and the omitted terms are non-negative.

    We now claim that the greedy method yields a solution that attains this upper-bound.
    Denote by $g_1 < g_2 < \dots < g_{t_g} \in [m]$ a solution produced by the greedy algorithm.
    Denote for convenience $g_0 = 0.$
    Let $\kappa \in [t_g].$ The candidate set in iteration $\kappa$ is
    \[
        \{\sigma \fw \gamma_{\sigma} < \gamma_{g_{\kappa-1}}\}.
    \]
    Since $\gamma$ values are non-increasing, it follows that this set equals
    \[
        \{\sigma \geq s_\kappa\}
    \]
    for some $s_\kappa > g_{\kappa-1}$.
    Hence, the choice in iteration $\kappa$, satisfies
    \[
        z_{g_\kappa} = u_{s_\kappa}.
    \]
    Now let $\sigma \in \{g_{\kappa-1} + 1, \dots, g_\kappa\}$ such that $\gamma_{\sigma} < \gamma_{\sigma-1} $.
    If that condition does not hold, then
    \[
        \gamma_{\sigma - 1} - \gamma_{\sigma} = 0,
    \]
    so $\sigma$ does not contribute to the bound derived in \eqref{eq:sep-pos-proof-1}.
    It follows that \[
                        \gamma_{\sigma} < \gamma_{\sigma-1} \leq \gamma_{g_{\kappa-1}}.
    \]
    So $\sigma \geq s_\kappa$.
    By definition of $u$, we then have that $u_{\sigma} \leq u_{s_\kappa} = z_{g_\kappa}.$
    Since $\sigma \leq g_\kappa$, by definition of $u$ it holds that $u_\sigma \geq z_{g_\kappa}.$
    So \[
           z_{g_\kappa} \leq u_{\sigma} \leq z_{g_\kappa}.
    \]
    Hence, equality holds on this entire range.

    For $\sigma > g_{t_g}, $ the greedy algorithm terminated, so there is no $\sigma'$ such that $\gamma_{\sigma'} < \gamma_{g_{t_g}}$,
    meaning $\gamma_\sigma = \gamma_{g_{t_g}}$ for all $\sigma > g_{t_g}$.
    Therefore
    \begin{equation}\label{eq:sep-pos-proof-2}
        \begin{split}
            S(\bm{g})
            & = \sum_{\kappa = 1}^{t_g} \sum_{\sigma = g_{\kappa-1}+1}^{g_\kappa}
                z_{g_\kappa}(\gamma_{\sigma-1} - \gamma_\sigma) \\
            & = \sum_{\kappa = 1}^{t_g} \sum_{\sigma = g_{\kappa-1}+1}^{g_\kappa}
                u_{\sigma}(\gamma_{\sigma-1} - \gamma_\sigma) \\
            & = \sum_{\sigma = 1}^m u_{\sigma}(\gamma_{\sigma-1} - \gamma_\sigma).
        \end{split}
    \end{equation}
    The second equality uses that $u_\sigma = z_{g_\kappa}$ whenever
    $\gamma_{\sigma-1} - \gamma_\sigma > 0$.
    The third uses that the ranges
    $\{g_{\kappa-1}+1, \dots, g_\kappa\}$ partition $\{1, \dots, g_{t_g}\}$
    and all remaining terms are zero.
    Combined with \eqref{eq:sep-pos-proof-1}, it follows that
    \[S(\bm{g}) \geq S(\bm{\tau})\]
    for any feasible sequence $\bm{\tau}$.
    So, $S(\bm{g})$ is a maximizer of \eqref{eq:sep-pos-sum}.
\end{proof}